\section{Introduction}
\label{sec:introduction}

Trust in political institutions is both an input into democratic stability and an output of how democracy is administered. Elections are the most visible point of contact between citizens and the state, so procedures that govern registration, identification, and vote verification can shape confidence in democratic institutions well beyond election day itself. These issues have become especially salient as electoral authorities around the world have modernized how citizens are registered and identified through biometric rolls and voter identification requirements, typically with the stated goal of strengthening electoral integrity and reducing fraud \citep{ideaBiometrics2017}. Yet these same reforms can also impose new administrative burdens, exclude citizens who fail to comply in time, and alter who remains politically represented \citep{ukECVoterID2024}. This paper studies that tension: when the state modernizes voter identification, does greater security translate into greater trust in democratic institutions, or does it instead narrow the electorate and invite political backlash?

I study this question through the staggered rollout of biometric voter registration (BVR) across Brazilian municipalities between 2008 and 2018. BVR required voters to register biometric information with the electoral authority and shifted polling-place verification away from the voter registration card and toward fingerprint authentication. In principle, this reform could improve registry accuracy, reduce impersonation and administrative error, and strengthen confidence in the electoral process. But it could also shrink the electorate if some citizens failed to complete the new registration procedure in time. I assemble a municipality-by-election panel from administrative electoral records and combine it with survey evidence from LAPOP to study this second wave of electoral modernization in Brazil, after the earlier transition to electronic voting at ballot casting \citep{Fujiwara2015VotingTechnoloPolitica}.

The administrative evidence shows that BVR reduced the size of the registered electorate at adoption, consistent with a combination of registry cleaning and incomplete compliance with the new registration requirements. The reform also changed the composition of the electorate, disproportionately reducing the share of lower-education voters on the rolls while increasing the relative share of more educated voters. At the same time, survey evidence suggests that exposure to this institutional modernization is associated with higher trust in institutions and in democracy, especially among lower-education respondents who remain in the electorate. These patterns suggest that biometric registration affected several margins at once: it made the voter registry more secure and arguably more accurate, but it also changed who was included in it. I also examine a distinct concern that has become more prominent in recent years---whether electoral technologies become focal points for distrust and partisan backlash. In the Brazilian context, I test whether Bolsonaro-aligned respondents in treated locations are systematically less trusting of vote-counting integrity or ballot secrecy, and I find little evidence that greater exposure to BVR amplified such distrust.

To organize these results, the paper introduces a compact formal model of electoral modernization before turning to the data. The model treats biometric registration as a security technology that raises the cost of remaining on the rolls, with those compliance costs borne more heavily by less-educated citizens. It yields four predictions. First, more demanding verification reduces registration disproportionately among low-education citizens. Second, because education and income are positively correlated, selective compliance shifts the post-registration median voter toward lower redistribution. Third, higher observed institutional trust among post-reform respondents can arise mechanically from selection into the surviving electorate rather than from persuasion, which is the key lens through which the LAPOP results should be read. Fourth, the planner's optimal biometric intensity is interior and depends on how much electoral integrity is valued relative to inclusion.

The paper's empirical leverage comes from the staggered municipal rollout of BVR, which creates variation in the timing of adoption across locations. I use event-study and difference-in-differences designs suited to staggered treatment timing to trace the dynamic effects of the reform on electorate size and composition, and I then pair those administrative results with survey evidence on institutional trust, democratic attitudes, and perceptions of electoral integrity. This setting allows the paper to make three contributions. First, it extends the literature on voting technology from ballot-casting technology to registration and voter-identification technology. Second, it contributes to the literature on administrative burden and political inclusion by showing that a security-oriented reform can improve the integrity of the registry while also excluding citizens who do not comply with new procedural requirements. Third, it speaks to political economy by showing that electoral modernization can shape not only the formal quality of election administration but also the composition of the electorate, citizens' trust in democratic institutions, and the scope for backlash against electoral technology. More broadly, the results suggest that the relationship between electoral security and democratic legitimacy runs through both the credibility of the process and the inclusiveness of the electorate.

\textbf{Literature overview.} \paragraph{Registration costs and the composition of participation.}
A first strand studies how electoral rules alter the cost of becoming, or remaining, an eligible voter. Classic evidence from the United States shows that registration requirements depress turnout, especially when deadlines are early, offices are inconvenient, or reregistration is required after moving \cite{Rosenstone1978TheEffectRegistra,Ansolabehere2005TheIntroducVoter}. Subsequent work complicates the distributional claim: education gaps need not always widen mechanically when registration rules change \cite{Nagler1991TheEffectRegistra}, while reforms that lower administrative frictions, such as preregistration and automatic reregistration, can raise participation among targeted groups \cite{Holbein2014MakingYoungVoters,Kim2022AutomatiVoterReregist}. Comparative work similarly emphasizes that turnout reflects both rules and mobilization incentives \cite{cox2015,frankMartinezComa2023,gaeblerEtAl2020}. The voter-ID literature pushes this argument toward institutional design that is explicitly justified by electoral integrity: strict ID rules can screen out voters without compliant documentation, with especially large burdens for racial and ethnic minorities and durable effects even when rules are later relaxed \cite{Fraga2021WhoVoterLaws,Hajnal2017VoterIdentifiLaws,Grimmer2021TheDurableDifferen}. Other evidence bounds the aggregate consequences more tightly, showing that in Florida and Michigan only a very small share of ballots were cast without ID and that even extreme disenfranchisement assumptions would rarely change election outcomes \cite{Hoekstra2019Strict}. This paper studies a related but distinct reform: biometric voter registration does not merely require identification at the polling place, but changes the administrative stock of registered voters itself, making it possible to measure both aggregate registry exit and compositional change.

\paragraph{Election technology and biometric state capacity.}
A second strand asks whether election technology expands effective access or instead substitutes one form of error for another. In Brazil, electronic voting reduced residual votes and increased the political weight of less-educated voters, with downstream effects on representation and infant health \cite{Fujiwara2015VotingTechnoloPolitica}. Yet the same broad technological modernization agenda can introduce new mistakes: electronic voting also increased party-label votes in ways consistent with ballot-interface confusion \cite{Zucco2016TradingOldErrors}. The closest existing study of Brazilian biometric registration finds that biometric recadastramento reduced the number of registered voters but did not significantly reduce actual turnout, interpreting the gap partly as the removal of invalid registrations \cite{forquesatoRodrigues2023}. Outside electoral administration, biometric monitoring has been studied as a state-capacity tool: biometric smartcards in India reduced leakage and improved payment delivery without clear evidence of beneficiary exclusion, while biometric monitoring in health care improved adherence and data accuracy \cite{Muralidharan2016BuildingStateCapacity,Bossuroy2024BiometriMonitoriService}. The unresolved issue is whether the integrity gains from biometric administration come with politically meaningful changes in who remains registered. This paper addresses that gap by decomposing registry change into exit and education relabeling at the municipality-cohort level.

\paragraph{Unequal turnout, representation, and welfare.}
A third branch links participation to representation and welfare. If turnout is uneven, policy need not reflect the preferences of the full eligible population: in city politics, higher turnout among underrepresented groups can shift descriptive and substantive representation \cite{Hajnal2005WhereTurnoutMatters}. Political-economy models give the mechanism theoretical bite: voting costs, candidate entry, and the location of the decisive voter shape equilibrium policy, redistribution, and the welfare consequences of participation rules \cite{besleyCoate1997,Meltzer1981RationalTheoryThe,Borgers2004CostlyVoting,Krasa2008MandatorVotingBetter}. A related public-finance literature shows why administrative restrictions can improve targeting in principle while also screening out real beneficiaries in practice \cite{Nichols1982TargetinTransferThrough,Deshpande2019WhoScreenedOut}. Recent evidence from Brazil reinforces the stakes: restrictions on fraudulent registration can alter political competition, candidate selection, public goods, and social outcomes \cite{Karim2025VoterBuyingPolitici}. The contribution here is to bring these welfare and representation concerns to a measured electoral-administration shock: rather than assuming that biometric recadastramento only removes ghosts or only burdens real voters, the paper estimates the size and composition of the affected electorate and uses those estimates in a welfare framework.

\paragraph{Electoral integrity, administration, and institutional trust.}
A fourth strand studies how citizens interpret electoral administration itself. Electoral reforms may increase procedural integrity, but they can also change beliefs about the fairness or competence of institutions. Mobilization around electoral procedures in Kenya increased turnout but lowered trust in the electoral commission, particularly among politically disappointed voters and in violent constituencies \cite{Marx2021VoterMobilisaAnd}. International and domestic actors also value electoral process independently of outcomes, which is why process integrity can become politically salient even when candidate stakes are high \cite{Bubeck2017ProcessCandidatThe}. Evidence from local election administration suggests that partisan administrators do not necessarily shift vote shares in favor of their party, but the question of perceived neutrality remains central to democratic legitimacy \cite{Ferrer2023HowPartisanLocal}. More broadly, electoral accountability can discipline corruption when citizens can observe and sanction incumbents \cite{ferrazFinan2011}. Biometric voter registration sits at the intersection of these concerns: it is advertised as an anti-fraud reform, but it may also make registration loss more visible, create confusion, or change who is represented in the electorate. The LAPOP component of this paper therefore treats trust in institutions and democratic attitudes as outcomes, not merely as background conditions.

\paragraph{Gap and contribution.}
Taken together, existing work shows that administrative frictions affect turnout, that election technologies can enfranchise some voters while generating new errors, that unequal participation matters for representation and welfare, and that procedural reforms can reshape trust. What remains less well understood is how a large biometric registration reform changes the electorate before votes are cast, especially outside the U.S. voter-ID setting and in a context where biometric registration is rolled out gradually across municipalities. This paper fills that gap by combining Brazilian administrative data with survey evidence: it estimates municipality-level registry exit, education relabeling, downstream voting patterns, and institutional trust, while distinguishing strict and hybrid implementation regimes. The result is a bridge between the voter-ID literature's concern with burdens, the state-capacity literature's focus on administrative integrity, and political-economy models of representation and welfare.


\textbf{Roadmap.} \hyperref[sec:background-data]{Section~\ref*{sec:background-data}} describes the institutional setting, the rollout of biometric registration, and the data. \hyperref[sec:empirical-design]{Section~\ref*{sec:empirical-design}} presents the empirical design. \hyperref[sec:electorate-size]{Section~\ref*{sec:electorate-size}} presents the main administrative evidence on the voter registry, covering both electorate size and composition. \hyperref[sec:decomposition]{Section~\ref*{sec:decomposition}} decomposes the registry shift into exit and education-record re-labeling. \hyperref[sec:downstream-outcomes]{Section~\ref*{sec:downstream-outcomes}} turns to political and downstream outcomes. \hyperref[sec:political-trust]{Section~\ref*{sec:political-trust}} turns to survey evidence on political trust, democratic attitudes, electoral-integrity beliefs, and backlash. \hyperref[sec:theory]{Section~\ref*{sec:theory}} develops the welfare framework, \hyperref[sec:calibration]{Section~\ref*{sec:calibration}} calibrates optimal electoral monitoring, and \hyperref[sec:conclusion]{Section~\ref*{sec:conclusion}} concludes.
